% Options for packages loaded elsewhere
\PassOptionsToPackage{unicode}{hyperref}
\PassOptionsToPackage{hyphens}{url}
%
\documentclass[
]{article}
\usepackage{amsmath,amssymb}
\usepackage{iftex,ctex}
\ifPDFTeX
  \usepackage[T1]{fontenc}
  \usepackage[utf8]{inputenc}
  \usepackage{textcomp} % provide euro and other symbols
\else % if luatex or xetex
  \usepackage{unicode-math} % this also loads fontspec
  \defaultfontfeatures{Scale=MatchLowercase}
  \defaultfontfeatures[\rmfamily]{Ligatures=TeX,Scale=1}
\fi
\usepackage{lmodern}
\ifPDFTeX\else
  % xetex/luatex font selection
\fi
% Use upquote if available, for straight quotes in verbatim environments
\IfFileExists{upquote.sty}{\usepackage{upquote}}{}
\IfFileExists{microtype.sty}{% use microtype if available
  \usepackage[]{microtype}
  \UseMicrotypeSet[protrusion]{basicmath} % disable protrusion for tt fonts
}{}
\makeatletter
\@ifundefined{KOMAClassName}{% if non-KOMA class
  \IfFileExists{parskip.sty}{%
    \usepackage{parskip}
  }{% else
    \setlength{\parindent}{0pt}
    \setlength{\parskip}{6pt plus 2pt minus 1pt}}
}{% if KOMA class
  \KOMAoptions{parskip=half}}
\makeatother
\usepackage{xcolor}
\usepackage{longtable,booktabs,array}
\usepackage{calc} % for calculating minipage widths
% Correct order of tables after \paragraph or \subparagraph
\usepackage{etoolbox}
\makeatletter
\patchcmd\longtable{\par}{\if@noskipsec\mbox{}\fi\par}{}{}
\makeatother
% Allow footnotes in longtable head/foot
\IfFileExists{footnotehyper.sty}{\usepackage{footnotehyper}}{\usepackage{footnote}}
\makesavenoteenv{longtable}
\setlength{\emergencystretch}{3em} % prevent overfull lines
\providecommand{\tightlist}{%
  \setlength{\itemsep}{0pt}\setlength{\parskip}{0pt}}
\setcounter{secnumdepth}{-\maxdimen} % remove section numbering
\ifLuaTeX
  \usepackage{selnolig}  % disable illegal ligatures
\fi
\IfFileExists{bookmark.sty}{\usepackage{bookmark}}{\usepackage{hyperref}}
\IfFileExists{xurl.sty}{\usepackage{xurl}}{} % add URL line breaks if available
\urlstyle{same}
\hypersetup{
  hidelinks,
  pdfcreator={LaTeX via pandoc}}

  \title{\Huge \textbf{马克思主义基本原理概论\\期末往年题重点}}
  \author{Zircon}
  \date{2023年秋季}
  \setcounter{tocdepth}{1}
  \usepackage{geometry}
  \geometry{left = 3cm, right = 3cm, top = 3cm, bottom = 3cm}
  
  \setcounter{secnumdepth}{3}


\begin{document}
\begin{titlepage}
  \vspace*{\fill}
  \begin{center}
    {\Huge\heiti \textbf{马克思主义基本原理概论\\期末往年题重点}}\\[30mm]
    {\Large Zircon}\\[5mm]
    2023年秋季
\end{center}
  \vspace*{\fill}
\end{titlepage}

\newpage
{\footnotesize \tableofcontents}

\hypertarget{ux7269ux8d28ux53caux5176ux5b58ux5728ux65b9ux5f0f}{%
\section{物质及其存在方式}\label{ux7269ux8d28ux53caux5176ux5b58ux5728ux65b9ux5f0f}}

\begin{quote}
\textbf{所谓物质，就是不依赖人的意识而存在，并能为人的意识所反映的客观存在。}马克思主义的物质范畴从现实存在着的自然与社会存在中抽象出了其共同特性------客观实在性。
\end{quote}

\begin{itemize}
\item
  根本属性：运动

  \begin{itemize}
  \item
    \textbf{物质的根本属性是运动。运动是标志着一切事物和现象的变化及其过程的哲学范畴。}
  \item
    物质和运动是不可分割的，运动是物质的运动，物质是运动着的物质，离开物质的运动和离开运动的物质都是不可想象的。
  \item
    运动的绝对性和相对性：

    \begin{itemize}
    \item
      物质世界的运动是绝对的，而物质在运动过程中又有某种相对的静止。
    \item
      相对静止是物质运动在一定条件下的稳定状态，具体包括两种状态：\\
      空间的相对位置的暂时不变 \& 事物的根本性质暂时不变。
    \item
      运动的绝对性体现了物质运动的变动性、无条件性，静止的相对性体现了物质运动的稳定性、有条件性。
    \item
      运动和静止相互依赖、相互渗透、相互包含，``动中有静、静中有动''。无条件的绝对运动和有条件的相对静止构成了对立统一的关系。
    \end{itemize}
  \end{itemize}
\item
  理论意义：

  \begin{itemize}
  \item
    \textbf{坚持了唯物主义一元论，同唯心主义一元论和二元论划清了界限}；
  \item
    \textbf{坚持了能动的反映论和可知论，批判了不可知论}；
  \item
    \textbf{体现了唯物论和辩证法的统一，克服了形而上学唯物主义的缺陷}；
  \item
    体现了\textbf{唯物主义自然观和唯物主义历史观的统一}，为彻底的唯物主义奠定了理论基础。
  \end{itemize}
\end{itemize}

\hypertarget{ux4e3bux89c2ux80fdux52a8ux6027ux548cux5ba2ux89c2ux89c4ux5f8bux6027ux7684ux8fa9ux8bc1ux7edfux4e00}{%
\section{主观能动性和客观规律性的辩证统一}\label{ux4e3bux89c2ux80fdux52a8ux6027ux548cux5ba2ux89c2ux89c4ux5f8bux6027ux7684ux8fa9ux8bc1ux7edfux4e00}}

\begin{itemize}
\item
  正确认识和把握物质和意识的辩证关系，还需要处理好主观能动性和客观规律性的关系。

  \begin{itemize}
  \item
    一方面，\textbf{尊重客观规律是正确发挥主观能动性的前提}。规律是事物变化发展过程中本身所固有的内在的、本质的、必然的联系。人们只有在认识和掌握客观规律的基础上，才能正确地认识世界，有效地改造世界。
  \item
    另一方面，\textbf{只有充分发挥主观能动性，才能正确认识和利用客观规律}。人能够通过自觉活动去认识规律，并按照客观规律去改造世界，以满足自身的需要。
  \item
    因此，尊重事物发展的客观规律性与发挥人的主观能动性是辩证统一的，\textbf{实践是客观规律性和主观能动性统一的基础}。
  \end{itemize}
\item
  正确发挥人的主观能动性，有以下三个方面的前提和条件：

  \begin{itemize}
  \item
    \textbf{从实际出发}是正确发挥人的主观能动性的前提；
  \item
    \textbf{实践}是正确发挥人的主观能动性的根本途径；
  \item
    正确发挥人的主观能动性，还需要依赖于一定的\textbf{物质条件和物质手段}。
  \end{itemize}
\end{itemize}

\hypertarget{ux5bf9ux7acbux7edfux4e00ux89c4ux5f8bux662fux552fux7269ux8fa9ux8bc1ux6cd5ux7684ux5b9eux8d28ux548cux6838ux5fc3}{%
\section{对立统一规律是唯物辩证法的实质和核心}\label{ux5bf9ux7acbux7edfux4e00ux89c4ux5f8bux662fux552fux7269ux8fa9ux8bc1ux6cd5ux7684ux5b9eux8d28ux548cux6838ux5fc3}}

对立统一规律是唯物辩证法的实质和核心。

\begin{itemize}
\item
  对立统一规律\textbf{揭示了事物普遍联系的根本内容和变化发展的内在动力}，从根本上回答了事物为什么会发展的问题；
\item
  对立统一规律是\textbf{贯穿量变质变规律、
  否定之否定规律以及唯物辩证法基本范畴的中心线索}，也是理解这些规律的``钥匙''；
\item
  对立统
  一规律提供了人们认识世界和改造世界的根本方法------\textbf{矛盾分析方法}。
\end{itemize}

\hypertarget{ux77dbux76feux5206ux6790ux6cd5}{%
\section{矛盾分析法}\label{ux77dbux76feux5206ux6790ux6cd5}}

\begin{quote}
\textbf{矛盾是反映事物内部和事物之间对立统一关系的哲学范畴}。对立和统一分别体现了矛盾的两种基本属性。矛盾的对立属性又称斗争性，矛盾的统一属性又称同一性。
\end{quote}

\begin{itemize}
\item
  矛盾的同一性和斗争性及其在事物发展中的作用：

  \begin{itemize}
  \item
    \textbf{同一性：矛盾着的对立面相互依存、相互贯通的性质和趋势}，有两个方面的含义：

    \begin{itemize}
    \item
      一是矛盾着的对立面\emph{相互依存}，互为存在的前提，并共处于一个统一体中；
    \item
      二是矛盾着的对立面相互贯通，在一定条件下可以\emph{相互转化}。
    \end{itemize}
  \item
    \textbf{斗争性：矛盾着的对立面相互排斥、相互分离的性质和趋势}，可以分为对抗性矛盾和非对抗性矛盾两种基本形式。
  \item
    关系：矛盾的同一性和斗争性相互联结、相辅相成。

    \begin{itemize}
    \item
      没有斗争性就没有同一性，没有同一性也没有斗争性，斗争性寓于同一性之中，同一性通过斗争性来体现；
    \item
      矛盾的同一性是有条件的、相对的，矛盾的斗争性是无条件的、绝对的；
    \item
      矛盾的同一性和斗争性相结合，构成了事物的矛盾运动，推动着事物的变化发展。
    \end{itemize}
  \item
    同一性在事物发展中的作用：

    \begin{itemize}
    \item
      \textbf{同一性是事物存在和发展的前提}；同一性使矛盾双方相互吸取有利于自身的因素，在相互作用中各自得到发展；
    \item
      \textbf{同一性规定着事物转化的可能和发展的趋势}。
    \end{itemize}
  \item
    斗争性在事物发展中的作用：

    \begin{itemize}
    \item
      矛盾双方的斗争促进矛盾双方力量的变化，造成双方力量发展的不平衡，为对立面的转化、事物的质变创造条件；
    \item
      \textbf{矛盾双方的斗争是一种矛盾统一体向另一种矛盾统一体过渡的决定性力量}。
    \end{itemize}
  \item
    在事物发展过程中，矛盾的同一性和斗争性相互结合，共同发生作用，但在不同条件下，二者所处的地位会有所不同。在一定的条件下，矛盾的斗争性可能处于主要方面，而在另外的条件下，矛盾的同一性又可能处于主要方面。
  \item
    矛盾同一性和斗争性的辩证关系原理，要求我们在观察和处理问题时，必须善于把两者结合起来，在斗争性中把握同一性，在同一性中把握斗争性。和谐是矛盾的一种特殊表现形式，体现着矛盾双方的相互依存、相互存进、共同发展。\textbf{和谐是相对的、有条件的，并不意味着矛盾的绝对同一}。
  \end{itemize}
\item
  矛盾的普遍性和特殊性及其相互关系：矛盾既有普遍性，又有特殊性。

  \begin{itemize}
  \item
    普遍性：矛盾存在于\textbf{一切事物}中，存在于\textbf{一切事物发展过程的始终}；
  \item
    特殊性：各个具体事物的矛盾、每一个矛盾的各个方面在发展的不同阶段上各有其特点。

    \begin{itemize}
    \item
      事物是由多种矛盾构成的。主要矛盾是矛盾体系处于支配地位、对事物发展起决定作用的矛盾。次要矛盾是矛盾体系中处于从属地位、对事物的发展起次要作用的矛盾。
    \item
      在每一对矛盾中，处于支配地位、起着主导作用的一方，是矛盾的主要方面；处于被支配地位、不起主导作用的一方，则是矛盾的次要方面。事物的性质是由主要矛盾的主要方面决定的。
    \item
      在实际工作中，要坚持``两点论''和``重点论''的统一：

      \begin{itemize}
      \item
        两点论是指在分析事物的矛盾时，既要看到矛盾双方的对立又要看到统一；既要看到主要矛盾、矛盾的主要方面，又要看到次要矛盾、矛盾的次要方面。
      \item
        重点论是指着重把握主要矛盾、矛盾的主要方面，并以此作为解决问题的出发点。
      \item
        ``两点论''和``重点论''的统一要求我们，看问题既要全面地看，又要看主流、大势、发展趋势。
      \end{itemize}
    \end{itemize}
  \item
    矛盾的普遍性和特殊性的辩证统一关系：

    \begin{itemize}
    \item
      矛盾的普遍性即矛盾的共性，矛盾的特殊性即矛盾的个性。
    \item
      矛盾的共性是无条件的、绝对的，矛盾的个性是有条件的、相对的。任何现实存在的事物的矛盾都是共性和个性的有机统一，共性寓于个性之中，没有离开个性的共性，也没有离开共性的个性。
    \end{itemize}
  \end{itemize}
\end{itemize}

\hypertarget{ux771fux7406ux7684ux7eddux5bf9ux6027ux548cux76f8ux5bf9ux6027}{%
\section{真理的绝对性和相对性}\label{ux771fux7406ux7684ux7eddux5bf9ux6027ux548cux76f8ux5bf9ux6027}}

\begin{quote}
真理是标志主观和客观相符合的哲学范畴，是对客观事物及其规律的正确反映。
\end{quote}

\begin{itemize}
\item
  真理的绝对性：真理的绝对性是指\textbf{真理主客观统一的确定性和发展的无限性}，具有两个方面的含义：

  \begin{itemize}
  \item
    任何真理都标志着\textbf{主观与客观相符合}，都包含着不依赖于人和人的意识的客观内容，都\textbf{同谬误有原则的界限}，这一点是绝对的、无条件的；
  \item
    \textbf{人类认识按其本性来说，能够正确认识无限发展着的物质世界}，认识每前进一步，都是对无限发展着的物质世界的接近，这一点也是绝对的、无条件的。
  \end{itemize}
\item
  真理的相对性：真理的相对性是指人们在一定条件下对客观事物及其本质和发展规律的正确认识总是\textbf{有限度的}、不完善的，也有两方面的含义：

  \begin{itemize}
  \item
    从客观世界的整体来看，任何真理都只是对客观世界某一阶段、某一部分的正确认识，人类已经达到的\textbf{认识的广度}总是有限度的\(\to\)认识有待扩展；
  \item
    就特定事物而言，任何真理都只是对客观对象一定方面、一定层次和一定程度的正确认识，\textbf{认识反映事物的深度}是有限度的，或是近似性的\(\to\)认识有待深化
  \end{itemize}
\item
  真理绝对性与相对性的辩证统一关系：

  \begin{itemize}
  \item
    二者\textbf{相互依存}：任何真理都既是绝对的，又是相对的；
  \item
    二者\textbf{相互包含}：绝对之中有相对，相对之中有绝对；
  \item
    \textbf{人类的认识永远处在由真理的相对性走向绝对性、接近绝对性的转化和发展过程中}。任何真理性的认识都是由真理的相对性向绝对性转化过程中的一个环节。真理既具有绝对性，又具有相对性，这是真理问题上的辩证法。
  \item
    \textbf{真理的绝对性与相对性根源于人类认识世界能力的无限性与有限性、绝对性与相对性的矛盾}。割裂真理的绝对性与相对性的辩证关系，就会走向形而上学的真理观，即绝对主义和相对主义。
  \end{itemize}
\end{itemize}

\hypertarget{ux771fux7406ux4e0eux8c2cux8bef}{%
\section{真理与谬误}\label{ux771fux7406ux4e0eux8c2cux8bef}}

\begin{quote}
\begin{itemize}
\item
  真理的定义：真理是标志着主观与客观相符合的哲学范畴，是对客观事物及其规律的正确反映。
\item
  谬误的定义：谬误是同客观事物及其发展规律相违背的认识，是对客观事物及其发展规律的歪曲反映。
\end{itemize}
\end{quote}

\begin{itemize}
\item
  真理与谬误的对立统一关系：真理与谬误是人类认识中的一对永恒矛盾，他们之间既对立又统一：

  \begin{itemize}
  \item
    真理与谬误相互对立。在确定的对象和范围内，其对立是绝对的，\textbf{存在着原则界限}。
  \item
    真理与谬误的对立又是相对的，在一定条件下能够相互转化：

    \begin{itemize}
    \item
      首先，\textbf{真理在一定条件下会转化为谬误}：

      \begin{itemize}
      \item
        一方面，\textbf{真理是具体的}。任何真理都是在一定范围内、一定条件下才能够成立，超出这个范围，失去特定条件，就会变成谬误；
      \item
        另一方面，\textbf{真理又是全面的}。把全面的真理性认识组成的科学体系中的某个原理孤立地抽取出来，切断同其他原理的联系，也会使其丧失自己的真理性而变为谬误。
      \end{itemize}
    \item
      其次，\textbf{谬误在一定条件下能够向真理转化}。
    \item
      最后，\textbf{在批判谬误中发展真理}，是谬误向真理转化的另一种形式。
    \end{itemize}
  \item
    真理和谬误的对立统一关系表明，真理总是同谬误相比较而存在、相斗争而发展的。
  \end{itemize}
\end{itemize}

\hypertarget{ux4e3bux89c2ux8fa9ux8bc1ux6cd5ux4e0eux5ba2ux89c2ux8fa9ux8bc1ux6cd5}{%
\section{主观辩证法与客观辩证法}\label{ux4e3bux89c2ux8fa9ux8bc1ux6cd5ux4e0eux5ba2ux89c2ux8fa9ux8bc1ux6cd5}}

唯物辩证法既包括客观辩证法也包括主观辩证法， 体现了唯物主义、
辩证法、认识论的统一。

\begin{quote}
\begin{itemize}
\item
  客观辩证法是指\textbf{客观事物或客观存在的辩证法}，
  即客观事物以相互作用、
  相互联系的形式呈现出的各种物质形态的辩证运动和发展规律。
\item
  主观辩证法是指\textbf{人类认识和思维运动的辩证法}，
  即以概念作为思维细胞的辩证思维运动和发展规律。
\end{itemize}
\end{quote}

\textbf{主观辩证法是客观辩证法在人的思维中的反映}，二者在本质上统一，但在表现形式上不同：

\begin{itemize}
\item
  客观辩证法采取\textbf{外部必然性形式}，不以人的意志为转移，是物质世界本身的联系和发展；
\item
  主观辩证法则采取\textbf{观念的、逻辑的形式}，是同人类思维的自觉活动相联系的，是以概念为基础的辩证思维规律，是辩证法的科学体系。
\end{itemize}

\hypertarget{ux5b9eux8df5ux51b3ux5b9aux8ba4ux8bc6}{%
\section{实践决定认识}\label{ux5b9eux8df5ux51b3ux5b9aux8ba4ux8bc6}}

\begin{itemize}
\item
  \textbf{实践是认识的来源}：

  \begin{itemize}
  \item
    认识的内容是在实践活动的基础上产生和发展的；
  \item
    离开实践的认识是不可能产生的。
  \item
    直接经验与间接经验：一切真知都是从直接经验发源的，但一个人的多数知识还是来自间接经验。
  \end{itemize}
\item
  \textbf{实践是认识发展的动力}：

  \begin{itemize}
  \item
    实践的需要推动认识的产生和发展；
  \item
    实践为认识的发展提供了手段和工具；
  \item
    实践改造了人的主观世界，锻炼和提高了人的认识能力。
  \end{itemize}
\item
  \textbf{实践是认识的目的}：人不是为了认识而认识，其最终目的是为实践服务，指导实践，以满足人们生活和生产的需要。
\item
  \textbf{实践是检验认识真理性的唯一标准}：认识是否具有真理性，只有在实践中才能得到验证。
\end{itemize}

\hypertarget{ux5b9eux8df5ux662fux68c0ux9a8cux771fux7406ux7684ux552fux4e00ux6807ux51c6}{%
\section{实践是检验真理的唯一标准}\label{ux5b9eux8df5ux662fux68c0ux9a8cux771fux7406ux7684ux552fux4e00ux6807ux51c6}}

实践之所以能够成为检验真理的唯一标准，是由真理的本性和实践的特点决定的。

\begin{itemize}
\item
  第一，从\textbf{真理的本性}看，真理是人们对客观事物及其发展规律的正确反映，它的本性在于主观和客观相符合。\textbf{检验真理就是检验人的主观认识同客观实际是否相符合以及符合的程度。}

  \begin{itemize}
  \item
    检验真理的标准，既不能是主观认识本身，也不能是客观事物。只有那种能够把主观认识与客观事物联系和沟通起来，从而使人们能够把二者加以比较和对照的东西，才能充当检验真理的标准。
  \item
    具有这种特性的东西，只能是作为主客观联系的桥梁的社会实践。
  \end{itemize}
\item
  第二，从\textbf{实践的特点}来看，\textbf{实践具有直接现实性}。实践的直接现实性是它的客观实在性的具体表现。

  \begin{itemize}
  \item
    如果实践的结果与实践之前的认识相符合，那么之前的认识就得到了证实，成为了真理性的证实。
  \item
    实践的直接现实性的品格，是实践能够成为检验真理唯一标准的主要根据。
  \end{itemize}
\item
  在实践检验真理的过程中，逻辑证明可以起到重要的补充作用。但是，\textbf{逻辑证明只能回答前提与结论的关系是否符合逻辑的问题，而不能回答结论是否符合客观实际的问题}。已被逻辑证明了的东西，还必须经过实践的检验，并最终服从实践检验的结果。
\end{itemize}

\hypertarget{ux65b0ux4e8bux7269ux4e0eux65e7ux4e8bux7269}{%
\section{新事物与旧事物}\label{ux65b0ux4e8bux7269ux4e0eux65e7ux4e8bux7269}}

\begin{quote}
发展的实质是\textbf{新事物的产生和旧事物的灭亡}。新事物是指合乎历史前进方向、具有远大前途的东西，旧事物是指丧失历史必然性，日趋灭亡的东西。
\end{quote}

新事物是不可战胜的：

\begin{itemize}
\item
  新事物与环境：新事物由于具有新的要素、结构和功能，适应已经变化了的环境和条件；
\item
  新事物与旧事物：事物是在旧事物的``母体''中孕育成熟的，对旧事物进行了``扬弃''，并添加了旧事物所不能容纳的新内容，使新事物在本质上优越于旧事物、具有强大生命力。
\end{itemize}

\hypertarget{ux793eux4f1aux5f62ux6001}{%
\section{社会形态}\label{ux793eux4f1aux5f62ux6001}}

\begin{quote}
\begin{itemize}
\item
  社会形态是关于社会运动的具体形式、发展阶段和不同质态的范畴，是同生产力发展一定阶段相适应的经济基础与上层建筑的统一体。
\item
  社会形态包括社会的经济形态、政治形态和意识形态，是三者的具体的、历史的统一。
\item
  一定的社会形态总要以一定的社会制度形式呈现出来，社会制度能够集中体现社会形态的性质，所以人们常用社会制度指代社会形态。
\item
  人类社会是不断发展的，社会的根本性变革和进步就是通过社会形态的更替实现的。
\end{itemize}
\end{quote}

\begin{itemize}
\item
  社会形态更替的统一性和多样性：

  \begin{itemize}
  \item
    依据生产关系的不同性质，社会历史可划分为五种社会形态：\textbf{原始社会、奴隶社会、封建社会、资本主义社会和共产主义社会}。这五种社会形态的依次更替，是\textbf{社会历史运动的一般过程和一般规律}，表现了社会形态更替的统一性。
  \item
    但是就某一国家或民族的社会发展历程而言，情况就复杂了。不同国家社会形态更替的过程不同；即便是同一种意识形态，在不同国家也会呈现出不同特点。所有这些都体现了社会形态更替形式的多样性。
  \end{itemize}
\item
  社会形态更替的必然性和人们的历史选择性：

  \begin{itemize}
  \item
    社会形态更替的统一性和多样性，根源于\textbf{社会发展的客观必然性与人们的历史选择性相统一}的过程。
  \item
    社会形态更替的客观必然性主要是指社会形态依次更替的过程和规律是客观的，其发展的基本趋势是确定不移的。\textbf{生产力与生产关系矛盾运动的规律性，从根本上规定了社会形态更替的客观必然性。}
  \item
    社会形态更替的规律也是人们自己的社会行动的规律。\textbf{规律的客观性并不否认人们历史活动的能动性}，并不排斥人们在遵循社会发展规律的基础上，对于某种社会形态的历史选择性：

    \begin{itemize}
    \item
      社会发展的客观必然性造成了一定历史阶段社会发展的基本趋势，\textbf{为人们的历史选择提供了基础、范围和可能性空间}；
    \item
      社会形态更替的过程也是一个\textbf{主观能动性和客观规律性相统一}的过程；
    \item
      \textbf{人们的历史选择性归根结底是人民群众的选择性}。人们对于社会形态的历史选择最终取决于人民群众的根本利益、根本意愿以及对社会发展规律的把握和顺应程度。
    \end{itemize}
  \end{itemize}
\item
  社会形态的更替的前进性与曲折性：

  \begin{itemize}
  \item
    社会形态的更替还表现为历史的前进性与曲折性、顺序性与跨越性的统一。
  \item
    社会形态更替的前进性、顺序性主要是指五种社会形态依次演进的基本趋势，其历史过程是一个``扬弃''的过程，但它并不否认历史发展的曲折性和跨越性。
  \item
    在社会进步发展的过程中，有时会出现社会形态更替的反复甚至倒退现象。
  \end{itemize}
\end{itemize}

\hypertarget{ux793eux4f1aux53d1ux5c55ux7684ux6839ux672cux52a8ux529b}{%
\section{社会发展的根本动力}\label{ux793eux4f1aux53d1ux5c55ux7684ux6839ux672cux52a8ux529b}}

\textbf{社会历史发展的根本动力是社会基本矛盾}。社会基本矛盾是指贯穿社会发展过程始终、规定社会发展过程的基本性质和基本趋势，并对社会历史发展起根本推动作用的矛盾。\textbf{生产力和生产关系、经济基础和上层建筑的矛盾是社会基本矛盾}。

\begin{itemize}
\item
  社会基本矛盾在历史发展中的作用主要表现在：

  \begin{itemize}
  \item
    首先，\textbf{生产力是社会基本矛盾运动中最基本的动力因素}，是人类社会发展和进步的最终决定力量；生产力是社会进步的根本内容，是衡量社会进步的根本尺度。
  \item
    其次，社会基本矛盾特别是生产力和生产关系的矛盾，\textbf{决定着社会中其他矛盾的存在和发展}；经济基础和上层建筑的矛盾也会影响和制约生产力和生产关系的矛盾。
  \item
    最后，\textbf{社会基本矛盾具有不同的表现形式和解决方式，并从根本上影响和促进社会形态的变化和发展}。
  \end{itemize}
\end{itemize}

\hypertarget{ux4ebaux6c11ux7fa4ux4f17ux662fux5386ux53f2ux7684ux521bux9020ux8005}{%
\section{人民群众是历史的创造者}\label{ux4ebaux6c11ux7fa4ux4f17ux662fux5386ux53f2ux7684ux521bux9020ux8005}}

\hypertarget{ux552fux7269ux53f2ux89c2ux8ba4ux4e3aux4ebaux6c11ux7fa4ux4f17ux662fux5386ux53f2ux7684ux521bux9020ux8005}{%
\subsection{唯物史观认为人民群众是历史的创造者}\label{ux552fux7269ux53f2ux89c2ux8ba4ux4e3aux4ebaux6c11ux7fa4ux4f17ux662fux5386ux53f2ux7684ux521bux9020ux8005}}

唯物史观与唯心史观的对立，在历史创造者的问题上表现为群众史观与英雄史观的对立。英雄史观从\textbf{社会意识决定社会存在}的基本前提出发，否认\textbf{物质资料的生产方式是社会发展的决定力量}，抹杀人民群众的历史作用，宣扬少数英雄人物创造历史。相反，群众史观认为历史的创造者不是个别英雄，而是人民群众。

\begin{itemize}
\item
  首先，唯物史观立足于\textbf{现实的人及其本质}来把握历史的创造者；
\item
  其次，唯物史观立足于\textbf{整体的社会历史过程}来探究谁是历史的创造者；
\item
  再次，唯物史观从\textbf{社会历史发展的必然性}入手考察和说明谁是历史的创造者；
\item
  最后，唯物史观从\textbf{人与历史关系的不同层次}上考察谁是历史的创造者。它不是对历史表象的经验描述，而是对历史本质的逻辑把握。
\end{itemize}

\hypertarget{ux4ebaux6c11ux7fa4ux4f17ux5728ux521bux9020ux5386ux53f2ux8fc7ux7a0bux4e2dux7684ux51b3ux5b9aux4f5cux7528}{%
\subsection{人民群众在创造历史过程中的决定作用}\label{ux4ebaux6c11ux7fa4ux4f17ux5728ux521bux9020ux5386ux53f2ux8fc7ux7a0bux4e2dux7684ux51b3ux5b9aux4f5cux7528}}

\begin{quote}
\begin{itemize}
\item
  从质上看，人民群众是指一切\emph{对社会历史发展起推动作用的人}；
\item
  从量上看，人民群众是指社会人口的绝大多数。
\item
  人民群众是一个历史范畴。在当代中国，全体社会主义劳动者、社会主义事业的建设者、拥护社会主义的爱国者、拥护祖国统一和致力于中华民族伟大复兴的爱国者都属于人民群众的范畴。
\end{itemize}
\end{quote}

在社会历史发展过程中，人民群众起着决定性的作用：

\begin{itemize}
\item
  人民群众是社会历史实践的主体，在创造历史中起决定性的作用；
\item
  人民群众是社会物质财富的创造者；
\item
  人民群众是社会精神财富的创造者；
\item
  人民群众是社会变革的决定力量。
\item
  人民群众的历史活动受到一定社会历史条件的制约。

  \begin{itemize}
  \item
    经济条件对于人民群众创造历史的活动有着首要的、决定性的影响。
  \end{itemize}
\end{itemize}

\hypertarget{ux65e0ux4ea7ux9636ux7ea7ux653fux515aux7684ux7fa4ux4f17ux8defux7ebf}{%
\subsection{无产阶级政党的群众路线}\label{ux65e0ux4ea7ux9636ux7ea7ux653fux515aux7684ux7fa4ux4f17ux8defux7ebf}}

\begin{itemize}
\item
  唯物史观关于人民群众是历史创造者的原理，要求我们坚持马克思主义群众观点，贯彻党的群众路线：

  \begin{itemize}
  \item
    群众观点：坚信人民群众自己解放自己的观点，一切向人民群众负责的观点，虚心向群众学习的观点。
  \item
    群众路线：群众路线是群众观点的具体应用，即一切为了群众、一切依靠群众，从群众中来，到群众中去。实质在于充分相信群众，坚决依靠群众，密切联系群众，全心全意为人民群众服务。
  \end{itemize}
\end{itemize}

\hypertarget{ux4e2aux4ebaux5728ux793eux4f1aux5386ux53f2ux4e2dux7684ux4f5cux7528}{%
\subsection{个人在社会历史中的作用}\label{ux4e2aux4ebaux5728ux793eux4f1aux5386ux53f2ux4e2dux7684ux4f5cux7528}}

任何历史人物的出现都体现了必然性与偶然性的统一。

\begin{itemize}
\item
  时势造英雄，杰出人物的出现具有必然性。
\item
  杰出人物对社会进程的产生影响仅仅是历史进程中的偶然现象，只能成为社会发展的个别原因。它们凭借自己的才能，虽然也能使具体历史事件的外貌或某些后果发生变化，但终究不能改变历史发展的基本方向。
\item
  如果看不到历史人物活动的社会制约性，割裂偶然与必然的关系，就势必会夸大个人的作用，进而否定或歪曲历史发展的规律。
\end{itemize}

\hypertarget{ux5546ux54c1ux4e8cux56e0ux7d20ux4e0eux52b3ux52a8ux4e8cux91cdux6027}{%
\section{商品二因素与劳动二重性}\label{ux5546ux54c1ux4e8cux56e0ux7d20ux4e0eux52b3ux52a8ux4e8cux91cdux6027}}

\begin{itemize}
\item
  商品二因素：\textbf{商品是用来交换、能满足人的某种需要的劳动产品}，具有使用价值和价值两个因素或两种属性，是\textbf{使用价值和价值的矛盾统一体}。

  \begin{itemize}
  \item
    使用价值：\textbf{商品能满足人的某种需要的有用性}，反映的是人与自然之间的物质关系，是商品的自然属性，是一切劳动产品所共有的属性，离开了它商品就不复存在。
  \item
    价值：\textbf{凝结在商品中的无差别的人类劳动，即人的脑力和体力的耗费}，价值是商品所特有的社会属性。
  \item
    使用价值和价值之间对立统一的关系：

    \begin{itemize}
    \item
      对立性：商品的使用价值和价值是相互排斥的，二者不可兼得。要获得商品的价值就必须放弃商品的使用价值；要得到商品的使用价值，就不能得到商品的价值。
    \item
      统一性：作为商品必须同时具有使用价值和价值两个因素。使用价值是价值的物质承担者，价值寓于使用价值之中。
    \end{itemize}
  \end{itemize}
\item
  劳动二重性：

  \begin{itemize}
  \item
    生产商品的劳动具有二重性，即\textbf{具体劳动和抽象劳动}。

    \begin{itemize}
    \item
      具体劳动：生产一定使用价值的具体形式的劳动。
    \item
      抽象劳动：撇开一切具体形式的、无差别的人类劳动，即人的脑力和体力的耗费。
    \end{itemize}
  \item
    \textbf{生产商品的具体劳动创造商品的使用价值，抽象劳动形成商品的价值}。任何一种劳动，一方面是特殊的具体劳动，另一方面又是一般的抽象劳动，这就是劳动的二重性。正是\textbf{劳动的二重性决定了商品的二因素}。
  \item
    具体劳动和抽象劳动也是对立统一的关系：

    \begin{itemize}
    \item
      一方面，具体劳动和抽象劳动不是各自独立存在的两种劳动或两次劳动，它们\textbf{在时间和空间上是统一的，是商品生产者的同一劳动过程不可分割的两个方面}；
    \item
      另一方面，具体劳动与抽象劳动又分别反映劳动的不同属性，\textbf{具体劳动所反映的是人与自然的关系，是劳动的自然属性，而抽象劳动所反映的是商品生产者的社会关系，是劳动的社会属性}。
    \end{itemize}
  \end{itemize}
\end{itemize}

\hypertarget{ux5546ux54c1ux7ecfux6d4eux7684ux57faux672cux77dbux76fe}{%
\section{商品经济的基本矛盾}\label{ux5546ux54c1ux7ecfux6d4eux7684ux57faux672cux77dbux76fe}}

\begin{quote}
\textbf{私人劳动和社会劳动的矛盾是商品经济的基本矛盾。}商品的使用价值和价值的矛盾、生产商品的具体劳动和抽象劳动的矛盾，根源于私人劳动和社会劳动的矛盾。
\end{quote}

\begin{itemize}
\item
  在以\textbf{私有制为基础的商品经济}中，\textbf{商品生产者的劳动既是具有社会性质的社会劳动，又是具有私人性质的私人劳动}：

  \begin{itemize}
  \item
    商品生产者的劳动的\textbf{社会性质是由社会分工决定的}。每个商品生产者的劳动都是社会总劳动的一部分，是具有社会性质的社会劳动。
  \item
    商品生产者的劳动的\textbf{私人性质是由生产资料私有制决定的}。商品生产者的劳动是按照自己的利益和要求进行的，是具有私人性质的私人劳动。
  \end{itemize}
\item
  私人劳动和社会劳动的矛盾构成私有制商品经济的基本矛盾，这一矛盾贯穿商品经济发展过程的始终，决定着商品经济的各种内在矛盾及其发展趋势：

  \begin{itemize}
  \item
    首先，私人劳动和社会劳动的矛盾决定着\textbf{商品经济的本质及发展过程}；

    \begin{itemize}
    \item
      商品经济是以交换为目的的经济形式，交换体现了商品经济的本质；
    \item
      交换是解决私人劳动和社会劳动之间矛盾的唯一途径。每个生产者的劳动本身是私人劳动，而私人劳动要转化为社会劳动，就必须用自己的产品去同别人的产品交换。
    \end{itemize}
  \item
    其次，私人劳动和社会劳动的矛盾是商品经济其他一切矛盾的基础；

    \begin{itemize}
    \item
      在商品交换过程中，\textbf{只有将具体劳动还原为抽象劳动，才能进行量的比较}。而具体劳动能否还原为抽象劳动，在根本上取决于\textbf{私人劳动和社会劳动能否实现统一}。
    \item
      商品生产者的私人劳动生产的产品如果与社会的需求不相适应，这种\textbf{私人劳动就不会被承认为社会劳动}，它作为\textbf{具体劳动的有用性质}也就不会为社会所承认，因而\textbf{不能还原为抽象劳动}，这意味着\textbf{商品的价值无法实现}，\textbf{商品的使用价值和价值之间的矛盾没有得到解决}。
    \end{itemize}
  \item
    最后，私人劳动和社会劳动的矛盾决定着商品生产者的命运。

    \begin{itemize}
    \item
      商品的售卖过程是私人劳动转化为社会劳动的过程。
    \end{itemize}
  \end{itemize}
\item
  私有制商品经济条件下，私人劳动和社会劳动之间的矛盾是商品经济的基本矛盾。在资本主义制度下，这种矛盾进一步发展成资本主义的基本矛盾，即\textbf{生产社会化和生产资料资本主义私人占有之间的矛盾}，正是这一矛盾的不断运动，使资本主义制度最终被社会主义制度所替代具有了客观必然性。
\end{itemize}

\hypertarget{ux5546ux54c1ux62dcux7269ux6559}{%
\section{商品拜物教}\label{ux5546ux54c1ux62dcux7269ux6559}}

\begin{quote}
私有制商品经济条件下\textbf{私人劳动与社会劳动之间的矛盾}通过商品的运动、
价值的运动、 货币的运动决定商品生产者的命运， 这使商品生产者认为商品、
价值乃至货币似乎是物的\textbf{自然属性}，
而这种所谓的自然属性又似乎具有一种\textbf{超自然的神秘性}，
商品生产者不能自己掌握自己的命运， 而是听凭商品、 价值、
货币运动的摆布，
\textbf{人与人之间一定的社会关系在人们面前采取了物与物的关系的虚幻形式}，
马克思称之为 ``商品拜物教＂ 。
\end{quote}

私有制商品经济条件下商品世界的拜物教性质的产生具有必然性。

\begin{itemize}
\item
  其一，\textbf{劳动产品只有采取商品的形式才能进行交换}，
  人类劳动的等同性只有采取同质的价值形式才能在交换中体现出来。
\item
  其二， \textbf{劳动量只有采取价值量这一物的形式才能进行计算和比较}。
\item
  其三，
  \textbf{生产者的劳动关系的社会性质只有采取商品之间即物与物之间相交换的形式才能间接地表现出来}，
  这就使人与人之间一定的社会关系被物与物的关系所掩盖，
  具有了拜物教性质。
  商品世界的拜物教性质或者生产商品的劳动的社会性质所具有的物的外观掩盖了商品经济关系的本质，
  妨碍人们透过物的外表认识商品、 价值以及货币的实质。
\end{itemize}

\hypertarget{ux4ef7ux503cux89c4ux5f8bux53caux5176ux4f5cux7528}{%
\section{价值规律及其作用}\label{ux4ef7ux503cux89c4ux5f8bux53caux5176ux4f5cux7528}}

\begin{quote}
价值规律是商品生产和商品交换的基本规律。价值规律贯穿于商品经济的全部过程，它既支配商品生产，又支配商品流通。
\end{quote}

\begin{itemize}
\item
  主要内容：\textbf{商品的价值量由生产商品的社会必要劳动时间决定，商品交换以价值量为基础，按照等价交换的原则进行。}
\item
  表现形式：\textbf{商品的价格围绕商品的价值自发波动。由于供求关系变化的影响，商品价格总是时而高于价值，时而低于价值，不可避免地围绕价值这个中心上下波动。}从较长时间来看，价格高于价值的部分和价格低于价值的部分能够相抵，商品的平均价格和价值是相一致的。
\item
  在市场配置资源过程中的作用：

  \begin{itemize}
  \item
    积极方面：自发地调节生产资料和劳动力在社会各生产部门之间的分配比例；自发地刺激社会生产力的发展；自发地调节社会收入分配。
  \item
    消极后果：导致资源浪费；阻碍技术进步；导致收入两极分化。
  \end{itemize}
\end{itemize}

\hypertarget{ux7b80ux5355ux52b3ux52a8ux4e0eux590dux6742ux52b3ux52a8}{%
\section{简单劳动与复杂劳动}\label{ux7b80ux5355ux52b3ux52a8ux4e0eux590dux6742ux52b3ux52a8}}

\begin{quote}
\begin{itemize}
\item
  简单劳动是指不需要经过专门训练和培养的一般劳动者都能从事的劳动。
\item
  复杂劳动是指需要经过专门训练和培养，
  具有一定文化知识和技术专长的劳动者所从事的劳动。
\end{itemize}
\end{quote}

商品的价值量同简单劳动与复杂劳动有密切的关系。
当复杂劳动生产出来的商品和简单劳动生产出来的商品相交换时，
商品的价值量是以简单劳动为尺度的。 复杂劳动等于自乘的或多倍的简单劳动，
也就是说，少量的复杂劳动等于多量的简单劳动。
\textbf{在相同的劳动时间里，
复杂劳动创造的价值大于简单劳动创造的价值}。在以私有制为基础的商品经济条件下，复杂劳动换算为简单劳动，不是商品生产者自觉计算出来的，
而是在\textbf{商品交换过程中自发实现的}。

\hypertarget{ux52b3ux52a8ux529bux6210ux4e3aux5546ux54c1ux4e0eux8d27ux5e01ux8f6cux5316ux4e3aux8d44ux672c}{%
\section{劳动力成为商品与货币转化为资本}\label{ux52b3ux52a8ux529bux6210ux4e3aux5546ux54c1ux4e0eux8d27ux5e01ux8f6cux5316ux4e3aux8d44ux672c}}

\begin{quote}
劳动力是指人的劳动能力，是人的脑力和体力的总和。劳动力的使用即劳动。
\end{quote}

\begin{itemize}
\item
  劳动力成为商品的基本条件：

  \begin{itemize}
  \item
    其一，劳动力是自由人，\textbf{能把自己的劳动力当作资本来支配}；
  \item
    其二，\textbf{劳动者没有任何生产资料，没有生活资料来源，因而不得不依靠出卖劳动力为生}。
  \end{itemize}
\item
  劳动力商品的特点：

  \begin{itemize}
  \item
    和任何商品一样，劳动力商品也有它的价值和使用价值。但是劳动力是特殊的商品，它的价值和使用价值具有不同于普通商品的特点。
  \item
    劳动力的价值是由生产、发展、维持和延续劳动力所必需的生活必需品的价值决定的，劳动力价值的最低界限是由生活上不可缺少的生活资料的价值决定的。一旦劳动力价值降低到这个界限以下，劳动力就只能在萎缩的状态下维持。
  \item
    （劳动力成为商品是货币成为资本的前提条件）劳动力商品在使用价值上有一个很突出的特点，就是它的\textbf{使用价值是价值的源泉}，在消费过程中能够创造新的价值，而且这个新的价值比劳动力本身的价值更大。正是由于这一特点，货币所有者购买到劳动力以后，在消费过程中，不仅能够收回他在购买这种特殊商品时支付的价值，还能得到一个增殖的价值即剩余价值(\(m\))。而\textbf{一旦货币购买的劳动力带来剩余价值，货币也就变成了资本}。
  \end{itemize}
\item
  资本总公式与劳动力成为商品

  \begin{itemize}
  \item
    资本首先表现为一定量的货币， 但货币本身并不就是资本。作为货币
    的货币和作为资本的货币，
    在形式上有明显区别。在``\(W-G-W\)"这个公式中， \(W\)代表商品，
    \(G\)代表流通中的货币， 这是商品流通公式，
    这个公式表明商品流通的目的是获得使用价值。在"\( G-W-G\)"这个公式中，
    \(W\)代表商品， \textbf{\(G\)代表的不是一般意义上的货币，
    而是作为资本的货币}，
    \(G'\)代表的是价值增殖后的货币。这个公式表明资本运动的一般目的是价值增殖，
    因此被称为 ``资本总公式'' 。
  \item
    从形式上看，
    \textbf{资本总公式与商品交换的原则是矛盾的}。价值规律要求商品交换遵循等价交换原则，交换领域不能创造新价值，
    但资本总公式却表明，
    资本在流通中创造了新价值。如何理解这个矛盾呢？问题的关键在于劳动力成为商品。（接着回答劳动力成为商品的基本条件）
  \end{itemize}
\end{itemize}

\hypertarget{ux65e0ux4ebaux5de5ux5382}{%
\section{``无人工厂''}\label{ux65e0ux4ebaux5de5ux5382}}

\begin{itemize}
\item
  无人工厂在本质上依然是物化劳动或不变资本的实物形式，它们在参加产品的生产时，只是把原有的价值转移到产品中去，而不创造新价值，更不能创造剩余价值。
\item
  在生产自动化的条件下，直接从事生产劳动的工人相对减少，而从事科研、设计、技术和管理劳动的人员日益增加，``总体工人''中脑力劳动者的比重不断增大。
\end{itemize}

\hypertarget{ux4e0dux53d8ux8d44ux672cux4e0eux53efux53d8ux8d44ux672c}{%
\section{不变资本与可变资本}\label{ux4e0dux53d8ux8d44ux672cux4e0eux53efux53d8ux8d44ux672c}}

资本在资本主义生产过程中采取生产资料和劳动力两种形态，根据这两部分资本在\textbf{剩余价值生产中所起的不同作用}，可将资本区分为不变资本和可变资本。

\begin{itemize}
\item
  不变资本(\(c\))：不变资本是以\textbf{生产资料}形态存在的资本。生产资料的价值通过工人的具体劳动被转移到新产品中，其转移的价值量不会大于它原有的价值量；
\item
  可变成本(\(v\))：可变成本是用来\textbf{购买劳动力}的那部分成本。可变资本的价值在生产过程中不是被转移到新产品上去，而是由工人的劳动再生产出来。在生产过程中，工人创造的新价值不仅包括相当于劳动力价值的价值，还包括一定量的剩余价值(\(m\))。
\end{itemize}

基于此，提出资本主义商品价值构成公式：\(W =c+v+m\).
将资本区分为不变资本和可变资本，进一步\textbf{解释了雇佣劳动者的剩余劳动是剩余价值的唯一源泉}。这种划分也为确定资本家对雇佣劳动者的剥削程度提供了科学依据。因此，要确定资本家对工人的剥削程度，就应该用剩余价值和雇佣劳动者的可变资本相比，即\(m' = \frac{m}{v}\)。

\hypertarget{ux5269ux4f59ux4ef7ux503cux751fux4ea7ux7684ux4e24ux79cdux57faux672cux65b9ux6cd5}{%
\section{剩余价值生产的两种基本方法}\label{ux5269ux4f59ux4ef7ux503cux751fux4ea7ux7684ux4e24ux79cdux57faux672cux65b9ux6cd5}}

\begin{itemize}
\item
  绝对剩余价值生产方法：在必要劳动时间不变的条件下，通过\textbf{延长工作时长和提高劳动强度}生产剩余价值。
\item
  相对剩余价值生产方法：通过\textbf{全社会劳动生产率的提高}实现。社会劳动生产率的提高，降低了劳动力的价值，从而缩短了必要劳动时间，相对延长了剩余劳动时间。（全社会劳动生产率的提高时资本家追逐超额剩余价值的结果。）
\end{itemize}

\hypertarget{ux8d44ux672cux79efux7d2f}{%
\section{资本积累}\label{ux8d44ux672cux79efux7d2f}}

\begin{quote}
把剩余价值转化为资本，或者说，\textbf{剩余价值的资本化}，就是资本积累。
\end{quote}

\begin{itemize}
\item
  本质：\textbf{资本家不断地利用无偿占有的工人创造的剩余价值来扩大自己的资本规模，进一步扩大和加强对工人的剥削与统治。}
\item
  一般规律：随着资本积累和生产规模的扩大，社会财富日益集中到资产阶级手中，而社会财富的直接创造者即无产阶级则只占有少部分社会财富。这样必然加剧社会的两极分化。
\item
  历史趋势：\textbf{资本主义制度的必然灭亡和社会主义制度的必然胜利}。

  \begin{itemize}
  \item
    一方面，资本主义生产越来越具有社会性；
  \item
    另一方面，生产社会化和生产资料资本主义私人占有之间的深刻矛盾日益加剧，资本主义最终会被新的、更能适应社会化大生产要求的社会形态所取代。
  \end{itemize}
\end{itemize}

\hypertarget{ux8d44ux672cux7684ux6709ux673aux6784ux6210}{%
\section{资本的有机构成}\label{ux8d44ux672cux7684ux6709ux673aux6784ux6210}}

\begin{quote}
\begin{itemize}
\item
  由生产的技术水平所决定的生产资料和劳动力之间的比例，叫做资本的技术构成；
\item
  不变资本和可变资本的价值比例，叫做资本的价值构成；
\item
  由资本的技术构成决定并反映技术构成变化的资本价值构成，叫做资本的有机构成。
\end{itemize}
\end{quote}

在资本主义生产过程中，资本有机构成的提高是一般趋势，这是由资本的本性决定的。

\begin{itemize}
\item
  资本主义生产的唯一动机和直接目的就是追求剩余价值。为了达到这一目的，资本家便尽可能改进技术以加快资本积累，通过资本积聚和资本集中扩大生产规模。
\item
  在自然形式上，每个劳动力所推动的生产资料的数量大幅增加；在价值形式上，不变资本部分日益增多，可变资本在资本总额中所占的比重日益下降，从而资本有机构成得以不断提高。
\item
  在资本有机构成提高的情况下，由于可变资本的相对量的减少，资本对劳动力的需求日益相对减少，其结果就是不可避免地造成大规模的工人失业，形成相对过剩人口。
\end{itemize}

\hypertarget{ux5de5ux8d44ux7684ux672cux8d28}{%
\section{工资的本质}\label{ux5de5ux8d44ux7684ux672cux8d28}}

\begin{quote}
在资本主义制度下，工人的工资表现为劳动力的价值或价格，这是资本主义工资的本质。
\end{quote}

\begin{itemize}
\item
  资本家购买工人的劳动力是以货币工资形式支付的，工人为资本家劳动，资本家付给工人工资，工资表现为劳动的价格或工人全部劳动的报酬，这就模糊了工人必要劳动和剩余劳动的界限，掩盖了资本主义的剥削关系。
\item
  在当代资本主义国家，随着社会生产力的发展、科学技术的进步、劳动过程的复杂化和脑力劳动作用的加强，工人的实际工资呈现出不断提高的趋势，但是与其创造的剩余价值的增长幅度相比，实际工资提高的幅度还是较小的。只要资本和雇佣劳动的基本经济关系不变，资本主义工资的本质就不会发生根本变化。
\end{itemize}

\hypertarget{ux5229ux6da6ux4e0eux5269ux4f59ux4ef7ux503c}{%
\section{利润与剩余价值}\label{ux5229ux6da6ux4e0eux5269ux4f59ux4ef7ux503c}}

\begin{itemize}
\item
  在现实的资本主义经济生活中，
  资本家并不是把剩余价值看作可变资本的产物，
  而是把它看作全部预付资本的产物或增加额， 剩余价值便取得了利润的形态。
\item
  当剩余价值转化为利润时，
  剩余价值率也转化为利润率。剩余价值率是剩余价值与可变资本的比率，
  反映的是资本家对工人的剥削程度；利润率则是剩余价值与全部预付资本的比率，
  反映的是预付总资本的增殖程度， 掩盖了资本家对工人的剥削。
\end{itemize}

\hypertarget{ux8d44ux672cux4e3bux4e49ux7684ux57faux672cux77dbux76fe}{%
\section{资本主义的基本矛盾}\label{ux8d44ux672cux4e3bux4e49ux7684ux57faux672cux77dbux76fe}}

\begin{quote}
\textbf{生产社会化和生产资料资本主义私人占有之间的矛盾}，是资本主义的基本矛盾。
\end{quote}

\begin{itemize}
\item
  资本主义基本矛盾的表现

  \begin{itemize}
  \item
    在资本主义条件下，随着科学技术的进步和社会生产力的不断提高，资本主义生产不断社会化。
  \item
    但是在资本家私人占有生产资料和剥削雇佣劳动者的生产关系中，\textbf{社会化的生产力却变成资本的生产力}，变成资本高效能地榨取剩余劳动、占有剩余价值、实现价值增殖的能力。
  \item
    这就形成了资本主义所特有的生产社会化和生产资料资本主义私人占有之间的矛盾。这是生产力和生产关系之间的矛盾在资本主义社会的具体体现。
  \end{itemize}
\end{itemize}

\hypertarget{ux8d44ux672cux4e3bux4e49ux7ecfux6d4eux5371ux673a}{%
\section{资本主义经济危机}\label{ux8d44ux672cux4e3bux4e49ux7ecfux6d4eux5371ux673a}}

\begin{quote}
资本主义发展到一定阶段，就会发生以\textbf{生产过剩}为基本特征的经济危机。
\end{quote}

\begin{itemize}
\item
  本质特征：生产过剩是资本主义经济危机的本质特征，但这种过剩是\textbf{相对过剩}，即\textbf{相对于劳动人民有支付能力的需求}来说社会生产的商品显得相对过剩，而不是与劳动人民的实际需要相比的绝对过剩。
\item
  形式原因：经济危机的抽象的一般的可能性，首先是由\textbf{货币作为流通手段和支付手段}引起的。

  \begin{itemize}
  \item
    以货币为媒介的商品买卖在时间上分为两个相互独立的行为。如果一些商品生产者在出卖自己的商品后不接着购买他人生产的商品，就会有另一些商品生产者的商品卖不出去；
  \item
    同时，在商品买卖有更多的部分采取赊购赊销方式的情况下，如果某些债务人在债务到期时不能支付，就会是整个信用关系遭到破坏；但是，这仅仅是危机形式上的可能性。
  \end{itemize}
\item
  根本原因：\textbf{资本主义的基本矛盾}，这一基本矛盾具体表现在以下两个方面：

  \begin{itemize}
  \item
    其一，\textbf{生产无限扩大的趋势与劳动人民有支付能力的需求相对缩小的矛盾}；
  \item
    其二，\textbf{单个企业内部生产的有组织性和整个社会生产的无政府状态之间的矛盾}。
  \end{itemize}
\item
  周期性：

  \begin{itemize}
  \item
    资本主义经济危机具有周期性，这是由\textbf{资本主义基本矛盾运动的阶段性}决定的。只要存在资本主义制度，经济危机就是不可避免的。
  \item
    资本主义经济危机周期性爆发的特点，使社会资本再生产也呈现出周期性的特点，从一次危机开始到另一次危机的爆发，就是再生产的一个周期，一般包括危机、萧条、复苏和高涨四个相互联系的阶段。其中，危机阶段是必经阶段，没有危机阶段，就不存在资本主义再生产的周期性。
  \end{itemize}
\end{itemize}

\hypertarget{ux5784ux65ad}{%
\section{垄断}\label{ux5784ux65ad}}

\hypertarget{ux5784ux65adux7684ux5f62ux6210}{%
\subsection{垄断的形成}\label{ux5784ux65adux7684ux5f62ux6210}}

\begin{quote}
垄断是指少数资本主义大企业为了获得高额利润，通过相互协议或联合，对一个或几个部门商品的生产、销售和价格进行操纵和控制。
\end{quote}

\begin{itemize}
\item
  产生的原因：

  \begin{itemize}
  \item
    当生产集中发展到相当高的程度，极少数企业就会联合起来，操纵和控制本部门的生产和销售，实行垄断，以获得高额利润；
  \item
    企业规模巨大，形成对竞争的限制，也会产生垄断；
  \item
    激烈的竞争给竞争各方带来的损失越来越严重，为了避免两败俱伤，企业之间会达成妥协，联合起来，实行垄断。
  \end{itemize}
\item
  垄断的方式：

  \begin{itemize}
  \item
    垄断是通过一定的垄断组织形式实现的。垄断组织是指在一个或几个经济部门中，占据垄断地位的大企业联合。
  \item
    尽管垄断组织的形式多种多样且不断变化发展，但本质上是一样的，即通过联合实现独占和瓜分商品生产和销售市场，操纵垄断价格，以攫取高额垄断利润。
  \end{itemize}
\end{itemize}

\hypertarget{ux5784ux65adux4e0bux7adeux4e89ux7684ux7279ux70b9}{%
\subsection{垄断下竞争的特点}\label{ux5784ux65adux4e0bux7adeux4e89ux7684ux7279ux70b9}}

垄断作为自由竞争的对立面，并不能消除竞争，反而使竞争变得更加复杂和激烈，这是因为：

\begin{itemize}
\item
  垄断没有消除产生竞争的经济条件；
\item
  垄断必须通过竞争来维持；
\item
  社会生产是复杂多样的，任何垄断组织都不可能把包罗万象的社会生产全部包下来；

  总之，在垄断条件下，在垄断组织内部、垄断组织之间以及垄断资本家集团之间，垄断组织同非垄断组织之间以及中小企业之间，存在广泛而激烈的竞争。
\end{itemize}

垄断条件下的竞争相比于自由竞争的新特点：\textbf{规模大、时间长、手段残酷、程度激烈，而且具有更大的破坏性}

\begin{longtable}[]{@{}lll@{}}
\toprule\noalign{}
& 自由竞争 & 垄断 \\
\midrule\noalign{}
\endhead
\bottomrule\noalign{}
\endlastfoot
竞争目的 & 获取更多的利润或超额利润，不断扩大资本的积累 &
\textbf{获取高额垄断利润}，并不断\textbf{巩固和扩大自己的垄断地位和统治权力} \\
竞争手段 &
主要运用经济手段，如通过改进技术、提高劳动生产率、降低产品成本 &
除了各种形式的经济手段之外，还采取\textbf{非经济的手段}，使竞争变得更加复杂激烈 \\
竞争范围 & 主要在经济领域、国内市场上进行 &
\textbf{国际市场上的竞争越来越激烈}，不仅经济领域的竞争多种多样，而且还扩大到经济领域之外 \\
\end{longtable}

\hypertarget{ux5784ux65adux5229ux6da6ux548cux5784ux65adux4ef7ux683c}{%
\subsection{垄断利润和垄断价格}\label{ux5784ux65adux5229ux6da6ux548cux5784ux65adux4ef7ux683c}}

\begin{quote}
垄断资本的实质在于获取垄断利润，垄断利润是垄断资本家凭借其在社会生产和流通中的垄断地位而获得的\textbf{超过平均利润}的高额利润。
\end{quote}

\begin{itemize}
\item
  垄断利润的来源：垄断资本所获得的高额利润，归根到底来\textbf{自无产阶级和其他劳动人民创造的剩余价值}。具体来说，垄断利润的来源大体有以下几个方面：

  \begin{itemize}
  \item
    来自对本国无产阶级和其他劳动人民剥削的加强；
  \item
    由于垄断资本可以通过垄断高价和垄断低价来控制市场，使得它能获得一些其他企业特别是非垄断企业的利润；
  \item
    通过加强对其他国家劳动人民的剥削和掠夺，从国外获取利润；
  \item
    通过资本主义国家政权进行有利于垄断资本的再分配，从而将劳动人民创造的国民收入的一部分变成垄断资本的收入。
  \end{itemize}
\end{itemize}

\begin{quote}
垄断利润主要是通过垄断组织指定的垄断价格来实现的。垄断价格是垄断组织在销售或购买商品时，凭借其垄断地位规定的、旨在保证获取最大限度利润的市场价格。
\end{quote}

\begin{itemize}
\item
  垄断价格的形式：垄断价格包括\textbf{垄断高价和垄断低价}两种形式。

  \begin{itemize}
  \item
    垄断高价是指垄断组织出售商品时规定的高于生产价格的价格；
  \item
    垄断低价是指垄断组织在购买非垄断企业所生产的原材料等生产资料时规定的低于生产价格的价格。
  \end{itemize}
\item
  垄断价格与价值规律：

  \begin{itemize}
  \item
    垄断组织操纵价格所带来的结果是抑制了市场上价格的自由波动，\textbf{垄断价格一定时期内背离生产价格和价值}。
  \item
    但是，从全社会看，\textbf{整个社会商品的价值仍然是由生产它们的社会必要劳动时间决定的}，垄断价格既不能增加也不能减少整个社会所生产的价值总量，它只是对商品价值和剩余价值作了\textbf{有利于垄断资本的再分配}。
  \item
    从全社会看，\textbf{商品的价格总额仍然等于商品的价值总额}。所以，垄断价格的产生没有否定价值规律，它是价值规律在垄断资本主义阶段作用的具体体现。
  \end{itemize}
\end{itemize}

\hypertarget{ux5784ux65adux8d44ux672cux4e3bux4e49ux76845ux4e2aux57faux672cux7279ux5f81}{%
\subsection{垄断资本主义的5个基本特征}\label{ux5784ux65adux8d44ux672cux4e3bux4e49ux76845ux4e2aux57faux672cux7279ux5f81}}

\begin{quote}
\begin{itemize}
\item
  金融资本：由\textbf{工业垄断资本和银行垄断资本}融合在一起而形成的一种垄断资本。
\item
  金融寡头：在金融资本形成的基础上产生，指\textbf{操纵国民经济命脉}，并在实际上\textbf{控制国家政权的少数垄断资本家}或垄断资本家集团。他们支配了大量的社会财富，控制了整个国家的经济命脉和上层建筑，是垄断资本主义国家事实上的统治者。
\end{itemize}
\end{quote}

\begin{itemize}
\item
  垄断组织在经济生活中起决定作用；
\item
  在金融资本的基础上形成金融寡头的统治；
\item
  资本输出有了特别重要的意义；
\item
  瓜分世界的资本家国际垄断同盟已经形成；
\item
  最大资本主义大国已把世界上的领土瓜分完毕。
\end{itemize}

\hypertarget{ux7ecfux6d4eux5168ux7403ux5316}{%
\section{经济全球化}\label{ux7ecfux6d4eux5168ux7403ux5316}}

\begin{quote}
经济全球化是指在生产不断发展、科技加速进步、社会分工和国际分工不断深化、生产的社会化和国际化程度不断提高的情况下，世界各国、各地区的经济活动越来越超出某一国家和地区的范围而相互联系、相互依赖的过程。
\end{quote}

\begin{itemize}
\item
  表现（进程）：\textbf{生产全球化；贸易全球化；金融全球化}
\item
  动因：

  \begin{itemize}
  \item
    从本质上来讲，经济全球化是\textbf{生产力发展和社会化大生产的必然要求}：
  \item
    \textbf{科学技术的进步和生产力的发展}为经济全球化提供了坚实的物质基础和根本的推动力；
  \item
    \textbf{跨国公司的发展}为经济全球化提供了适宜的企业组织形式；
  \item
    \textbf{各国经济体制的变革和国际经济组织的发展}是经济全球化的体制与组织保障。
  \end{itemize}
\item
  影响：

  \begin{itemize}
  \item
    经济全球化的过程是\textbf{生产社会化程度不断提高}的过程。在全球化进程中，社会分工得以在更大的范围内进行，\textbf{资金、技术等生产要素可以在国际社会流动和优化配置}，由此带来巨大的分工利益，推动世界生产力的发展。
  \item
    一些发展中国家在参与经济全球化的进程中也得到了快速发展：

    \begin{itemize}
    \item
      经济全球化为发展中国家提供\textbf{先进技术和管理经验}；
    \item
      经济全球化为发展中国家提供更多的\textbf{就业机会}；
    \item
      经济全球化推动发展中国家\textbf{国际贸易发展}；
    \item
      经济全球化促进发展中国家\textbf{跨国公司的发展}。
    \end{itemize}
  \item
    经济全球化也是一把``双刃剑''：

    \begin{itemize}
    \item
      \textbf{发达国家和发展中国家在经济全球化过程中的地位和收益不平等、不平衡}；
    \item
      加剧了发展中国家\textbf{资源短缺和环境污染}；
    \item
      一定程度上增加\textbf{经济风险}。
    \end{itemize}
  \end{itemize}
\end{itemize}

\hypertarget{ux8d44ux672cux4e3bux4e49ux7684ux5386ux53f2ux5730ux4f4dux548cux53d1ux5c55ux8d8bux52bf}{%
\section{资本主义的历史地位和发展趋势}\label{ux8d44ux672cux4e3bux4e49ux7684ux5386ux53f2ux5730ux4f4dux548cux53d1ux5c55ux8d8bux52bf}}

\hypertarget{ux8d44ux672cux4e3bux4e49ux7684ux5386ux53f2ux5730ux4f4d}{%
\subsection{资本主义的历史地位}\label{ux8d44ux672cux4e3bux4e49ux7684ux5386ux53f2ux5730ux4f4d}}

\begin{itemize}
\item
  与封建社会相比，资本主义显示了巨大的\textbf{历史进步性}：

  \begin{itemize}
  \item
    \textbf{资本主义将科学技术转变为强大的生产力}；
  \item
    资本主义\textbf{追求剩余价值的内在动力和竞争的外在压力推动了社会生产力的迅速发展}；
  \item
    资本主义的\textbf{意识形态和政治制度}作为上层建筑在战胜封建社会自给自足的小生产的生产方式，保护、促进和完善资本主义生产方式方面起着重要作用，从而推动了社会生产力的迅速发展，促进了社会进步。
  \end{itemize}
\item
  资本主义的局限性：

  \begin{itemize}
  \item
    \textbf{资本主义基本矛盾阻碍社会生产力的发展}；
  \item
    资本主义制度下\textbf{财富占有两极分化，引起经济危机}；
  \item
    资本家阶级支配和控制资本主义经济和政治的发展和运行，\textbf{不断激化社会矛盾和冲突}。
  \end{itemize}
\end{itemize}

\hypertarget{ux8d44ux672cux4e3bux4e49ux88abux793eux4f1aux4e3bux4e49ux6240ux66ffux4ee3ux7684ux5386ux53f2ux5fc5ux7136ux6027}{%
\subsection{资本主义被社会主义所替代的历史必然性}\label{ux8d44ux672cux4e3bux4e49ux88abux793eux4f1aux4e3bux4e49ux6240ux66ffux4ee3ux7684ux5386ux53f2ux5fc5ux7136ux6027}}

资本主义的内在矛盾决定了资本主义必然被社会主义所替代。

\begin{itemize}
\item
  \textbf{资本主义基本矛盾包含着现代的一切冲突的萌芽}：

  \begin{itemize}
  \item
    资本主义基本矛盾表现在\textbf{阶级关系上，是无产阶级和资产阶级的对立}；
  \item
    资本主义基本矛盾表现在\textbf{生产上，是个别企业内部生产的有组织性和整个社会生产的无政府状态之间的对立}。
  \item
    资本主义\textbf{经济危机的爆发}正是这个基本矛盾发展的结果。只有用社会主义生产方式取而代之，才能根本解决资本主义生产方式的基本矛盾。
  \end{itemize}
\item
  \textbf{资本积累推动资本主义基本矛盾不断激化并最终否定资本主义自身}：

  \begin{itemize}
  \item
    \textbf{资本主义基本矛盾在基本积累过程中不断发展，资本的不断积累为否定资本主义制度自身准备了物质条件}。
  \end{itemize}
\item
  \textbf{国家垄断资本主义是资本社会化的更高形式，将成为社会主义的前奏}：

  \begin{itemize}
  \item
    资本国有化将为社会主义革命提供直接的物质前提，是无产阶级社会主义革命的入口处。
  \end{itemize}
\item
  \textbf{资本主义社会存在和资产阶级和无产阶级两大阶级之间的矛盾和斗争}：

  \begin{itemize}
  \item
    随着资本主义经济的发展，资产阶级由生产力的解放者日益变成阻碍者；
  \item
    资本主义在造就了社会化大生产的同时，也产生了推动和运用这一先进生产力的无产阶级，在经济上所处的被剥削地位使无产阶级具有彻底的革命性和斗争精神。\textbf{社会化大生产使无产阶级成为最有组织性的革命力量}。
  \end{itemize}
\item
  当今世界，虽然\textbf{资本主义制度通过自我调节还能为生产力的发展提供一定的空间}，但社会主义取代资本主义是历史的必然趋势。

  \begin{itemize}
  \item
    马克思、恩格斯关于资本主义社会基本矛盾的分析没有过时，关于资本主义必然消亡、社会主义必然胜利的历史唯物主义观点也没有过时，这是社会历史发展不可逆转的总趋势；
  \item
    但道路是曲折的。资本主义最终消亡、社会主义最终胜利，必然是一个很长的历史过程。我们要深刻认识资本主义社会的自我调节能力，充分估计到西方发达国家在经济科技军事方面长期占据优势的客观现实，认真做好两种社会制度长期合作和斗争的各方面准备。
  \end{itemize}
\end{itemize}

\end{document}
